\section{Discussion}\label{sec:discussion}

\subsection{Map-level inference at survey scale}\label{sec:discussion:maplevel}

Because the map-level posterior is trained on simulations that run the cluster finder on multi-frequency
maps, the effects of the finder enter its forward model by construction. These include the confusion
noise, the search lift, neighbours and the full shape of the finder's errors, upward tail included, so that
we need neither a functional form for these errors nor fits at chosen points of the prior. In our tests,
this approach gives posteriors within 0.12 widths of the truth and with close to correct coverage at both
survey sizes (Section~\ref{sec:results}). We now discuss what it costs, what it gains as the survey grows
and what still limits it.

\paragraph*{Cost.} At 40.7 per cent of the sky, which is about the area of SO \citep{SO2019}, one survey
simulation takes 1.6\,s on one GPU and a training set of 10\,000 surveys takes 4.6\,h
(Table~\ref{tab:cost}). With the present simulator, map-level SBI is therefore affordable at the scale of
SO. We note that testing any posterior, whether map-level or catalogue-level, requires map-level
simulations as well, since the bias and coverage tests of Section~\ref{sec:results} use 100 map-level mocks
and 500 map-level prior draws.

\paragraph*{Gain with survey area.} The value of a correct forward model can grow with the size of the
survey. If an error in the forward model shifts the posterior mean of a parameter $\theta$ by $\Delta$, the
bias statistic of equation~(\ref{eq:bias}) becomes
\begin{equation}
\mathcal{B} \simeq \frac{\Delta}{\mathrm{sd}[\theta\mid\bm x]},
\label{eq:disc_bias}
\end{equation}
where $\Delta$ depends on the model error but hardly on the survey area, whereas the posterior width
shrinks as the area grows. For a parameter constrained by the counts alone, the width scales as
$f_\mathrm{sky}^{-1/2}$, with $f_\mathrm{sky}$ the fraction of the sky covered, so that
$\mathcal{B}\propto f_\mathrm{sky}^{1/2}$ grows by a factor of 2 from 10 to 40.7 per cent of the sky. For
the uncalibrated reference, whose shift is large, the bias in widths grows by about 1.5 instead, because
the shift of its posterior mean hardly changes with area while its widths shrink 1.3--1.7 times
(Section~\ref{sec:results:bias}). We expect the widths to shrink by less than a factor of 2 for two reasons,
although we have not tested either. Six of the ten parameters have Gaussian priors that the counts do not
tighten, so that their uncertainty propagates into the four main parameters, and the edges of the flat
prior may also limit the shrinking. The calibrated models, in contrast, do not show this growth, since
their bias in $\alsz$ changes little between the two survey sizes, from 0.50 to 0.52 widths with the
one-point calibration and from 0.62 to 0.67 widths with the noise-level calibration
(Table~\ref{tab:bias}). Their shift thus does not stay fixed as the area grows, for reasons we have not
determined. In general, however, a model error whose shift does not change with area can matter more at
the area of SO than at 10 per cent of the sky. The map-level posterior, in contrast, stays within 0.12 widths at both
sizes.

\paragraph*{Current limits.} The map-level approach is limited, first, by the simulator, which needs the
additions of Section~\ref{sec:discussion:limits} before it is applied to data, each of which makes a
simulation more expensive. Second, the neural posterior at 10 per cent of the sky gives $\alsz$ intervals that are slightly
too narrow (90 per cent coverage of 0.86), and it may have a small offset in $\omc$
(Section~\ref{sec:discussion:omegac}). Finally, with the 21-element data vector of \citet{Zubeldia2025} as
its summary, the map-level posterior is not more precise than the catalogue-level ones
(Fig.~\ref{fig:corner}), although a richer summary can be added without changing the forward model
(Section~\ref{sec:discussion:limits}).

\subsection{Catalogue models calibrated on maps}\label{sec:discussion:calibrated}

Analyses of real data usually estimate the noise of the filter on the observed maps \citep[e.g.][]{Planck2016XXIV,Bocquet2019,Hilton2021}, and
this noise includes the confusion noise of the observed sky. Of the models in this paper, our one-point
calibration comes closest to this practice, because it uses the finder's noise measured on map-level mocks
at the true parameters and also fits the search lift, the scatter and the losses there
(Section~\ref{sec:calib-point}). At the true parameters it is within 0.2 widths of the truth in $\sigma_8$,
$\omc$ and $\asz$, while $\alsz$ is 0.50 widths high at 10 per cent of the sky and 0.52 widths high at 40.7
per cent (Table~\ref{tab:bias}). Halo by halo, the scale factor is $s=0.990$ against the finder's own noise
but 0.921 against the noise of the instrument and the CMB (Section~\ref{sec:finder_fit}), which shows that
most of the offset of the uncalibrated model comes from the noise. A catalogue model that used the measured
noise of each test sky with unit Gaussian scatter would separate the effect of the noise from those of the
lift, the scatter and the tail, but we have not run such a model.

\paragraph*{Shortcomings of the calibration.} A calibration has to carry the three effects of the finder
and the upward tail. The catalogue model leaves out three effects of the finder: confusion noise, the search lift
and neighbours; and it assumes Gaussian scatter, which the finder's errors break through an upward tail.
After masking, the confusion noise makes the finder's noise 4--8 per cent higher than that of the instrument and the CMB. All of these effects change with the amount of tSZ signal on the sky, and hence
with the parameters (Section~\ref{sec:calib-grid}). A calibration at a single point does not follow these
changes. Across the prior, the 90 per cent intervals of the one-point calibration therefore contain the true
$\sigma_8$ for only 0.65 of the skies at 10 per cent of the sky and 0.59 at 40.7 per cent
(Table~\ref{tab:coverage}). A calibration that follows the parameters through the noise level, fitted at
25 points, brings $\sigma_8$, $\omc$ and $\asz$ close to correct coverage, with 0.86--0.90 at the 90 per
cent level. However, it does not follow the slope, since $\alsz$ stays 0.62 and 0.67 widths high at the
true parameters and its 90 per cent coverage is 0.78 and 0.74, at 10 and 40.7 per cent of the sky,
respectively.

\paragraph*{Requirements for a better calibration.} We attribute the remaining bias in $\alsz$ to the
upward tail of Section~\ref{sec:finder_tail}, which changes with $q$ and with the amount of tSZ signal on the
sky. Our functional form, a Gaussian scatter around $\sqrt{(sq)^2+k}$ with a loss rate
(equation~\ref{eq:qobs_fit}), cannot carry such a tail. Splitting the calibration by cluster size does not change
this bias (Section~\ref{sec:calib-size}). To carry the tail, a catalogue model would need to emulate the
full error distribution of the finder as a function of $q$ and of the parameters, and building and
validating such an emulator also costs map-level simulations. Our calibrations use 25 points with 5 mocks
each, or 125 map-level survey simulations, compared with the 10\,000 of a map-level training set. The tests of
any catalogue model, however, need map-level mocks and prior draws as well
(Section~\ref{sec:discussion:maplevel}), so that once a map-level simulator exists and has been validated,
training on it adds only the cost of one training set.

None of these effects is specific to our finder, since any matched filter that estimates its noise from the
maps sees the tSZ signal of faint haloes, any search over positions and sizes has a search lift, and any
finder meets neighbours. How much these effects matter depends on the depth of the survey, on the noise
used by the catalogue model and on the priors (Section~\ref{sec:discussion:scope}), so that other surveys
and noise levels would need their own simulations.

\subsection{Scope of the uncalibrated reference}\label{sec:discussion:scope}

We include the uncalibrated catalogue model, which, as in forecasts, uses a noise table for the instrument
and the CMB only, as a reference and not as a test of analyses that calibrate on maps. Since both
forward models share the haloes, the scaling relations, the instrument and the CMB
(Section~\ref{sec:sims}), its bias comes from the observation step, that is, from the finder and the
confirmation. On our skies, the uncalibrated model is biased by up to 5.3 widths at 10 per cent of the sky
and 8.3 widths at 40.7 per cent, mostly in $\sigma_8$ and $\alsz$. In $S_8$, by contrast, it is 0.30 widths
low at 10 per cent of the sky (not measured at 40.7 per cent). This bias therefore measures the size of the
finder's effects at goal depth, and it depends on three choices, which we discuss in turn.

\paragraph*{Noise level.} We adopt the goal-depth forecast noise of \citet{SO2019}, the deepest of its
levels. At lower noise, the finder detects more faint clusters, for which the search lift matters most,
and the confusion noise makes up a larger part of the total. Conversely, at higher instrument noise the
confusion noise, which adds in quadrature, makes up a smaller part. We therefore expect a smaller bias at
higher noise, and closer agreement with the finder's counts than at goal depth, where the counts above
$q=8.4$ are 9--25 per cent below the model at the true parameters (Section~\ref{sec:finder_counts}),
although we have not run other noise levels.

\paragraph*{Noise of the catalogue model.} A noise measured on the maps includes the confusion noise.
With it, and with the fitted lift, scatter and losses of our one-point calibration, most of the bias
disappears at the true parameters, leaving a slope bias of about half a width and a coverage that depends
on how the calibration follows the parameters (Section~\ref{sec:discussion:calibrated}).

\paragraph*{Prior on the slope.} The prior on $\alsz$ is flat, $\mathcal U(1.00,1.24)$. Most of the bias goes into the slope and its degeneracy with $\sigma_8$, as the posterior explains the deficit of
bright clusters on the maps (Fig.~\ref{fig:counts}) with a shallower tSZ--mass relation. On the mock of
Fig.~\ref{fig:corner40}, the uncalibrated posterior lies about three widths above the lower edge of the
prior, so that the edge truncates little of it. Informative priors on the tSZ--mass relation, for example
from weak-lensing mass calibration, would restrict this direction. Although we have not tested this, we
expect them to reduce the bias, with part of the error possibly going into the cosmological parameters
instead.

\subsection{Limitations}\label{sec:discussion:limits}

Every test in this paper compares inferences on the same simulated skies, so that our conclusions rest on
the difference between the two forward models rather than on how realistic the skies are. An analysis of
real data, however, would need a more realistic simulator, and we discuss each of our simplifications in
turn.

\paragraph*{No halo clustering.} We place the haloes at uniform random positions on each tile and
therefore neglect halo clustering. On the real sky, massive clusters sit in overdense regions and are
surrounded by more haloes than average. The tSZ signal projected onto a cluster by other haloes is then
correlated with the cluster's mass, and the confusion noise around a massive cluster is larger than its
average over the tile. We therefore expect the upward tail of the finder's errors
(Section~\ref{sec:finder_tail}) and the effects of neighbours to be larger on the real sky than in our
simulations, although we have not simulated this. In addition, clustering adds sample variance to the counts,
which would widen the posteriors of both forward models. For these reasons, the simulator needs clustered
haloes before it is applied to data.

\paragraph*{Gaussian stationary noise.} Our noise is Gaussian, with the same power spectrum everywhere on a
tile and on every tile (equation~\ref{eq:noise}), whereas real maps have a depth that varies across the
footprint, scan-synchronous stripes and patchy atmosphere. Because the finder estimates its covariance per
tile, it can follow slow changes in depth but not changes within a tile. A real analysis would simulate the
noise from the data or from simulations of the scanning, which the map-level approach can include
directly.

\paragraph*{No foregrounds.} The maps contain neither the cosmic infrared background (CIB), which matters
most at high frequencies, nor radio sources, which matter most at low frequencies. Both are correlated
with clusters, since radio galaxies in clusters fill in part of the tSZ decrement at low frequency
\citep{Gupta2017} and the CIB adds emission where the tSZ signal is an excess \citep{Melin2018}.
In a standard matched filter, the weights of the channels at each multipole, $C(\ell)^{-1}\bm a(\ell)$,
minimise the variance of $\hat y_0$ under the single condition that the response to the tSZ spectrum is
unity. Emission that is correlated with the cluster, such as the CIB of its member galaxies, is therefore
not suppressed as noise, and part of it passes into $\hat y_0$. In Planck-like mocks from the Websky
simulation, \citet{Zubeldia2023} found that this biases the tSZ observables by about half a standard
deviation per cluster, and leads to about 20 per cent fewer detections when all the Planck high-frequency
channels are used. Their spectrally constrained filter, which is implemented in \textsc{SZiFi}, adds a
second condition, namely a zero response to an assumed CIB spectrum $\bm b(\ell)$, a modified blackbody.
Optionally, it also sets a zero response to the derivatives of that spectrum with respect to its spectral
index and temperature, which makes the filter robust to errors in the assumed spectrum. With a single
constraint, the weights become proportional to $C^{-1}\bm a-(\bm a^{\top}C^{-1}\bm b/\bm b^{\top}C^{-1}\bm
b)\,C^{-1}\bm b$. The CIB correlated with the cluster is thus removed, at the cost of a larger
$\sigma_\theta$ and hence a lower $q$. In our finder, this change would affect only the weights of each
bin in $\ell$. In a map-level simulator that paints the CIB, however, the bias that it causes would
already be part of the forward model. The constrained filter would then serve mainly to make the
inference less sensitive to how the CIB is modelled, at some cost in signal-to-noise. Which of the two
filters gives the more precise and more robust posterior could be measured with the test used here;
we have not simulated foregrounds. As with clustering, a catalogue model would have to describe foregrounds correlated with clusters separately,
whereas a map-level simulator includes them once they are painted.

\paragraph*{Idealised confirmation.} We confirm a peak if a halo above $5\times10^{13}\,\msun$ lies within
$1.5\arcmin$ of it, and assign to it that halo's exact redshift. Real catalogues, in contrast, are
confirmed by optical and infrared follow-up, with its own completeness, redshift errors and false matches.
In our simulations, noise peaks are removed unless a massive halo lies within $1.5\arcmin$ by chance (Section~\ref{sec:finder_zoom}). In addition, the redshift comes from the halo with the largest $y_0$ near the peak,
which need not be the halo that produced it. Since the catalogue model has neither effect, a real analysis
needs a model of the confirmation in either forward model.

\paragraph*{Flat periodic tiles.} We model the survey as a set of flat, periodic $10\degr\times10\degr$
tiles without masks or survey edges. Since the filter sizes, up to $15\arcmin$, are small compared with the
tile, the flat-sky approximation is adequate for the finder. Masks around point sources and survey edges,
however, change the noise of the filter locally and would need to be simulated.

\paragraph*{One pressure profile and a fixed mass bias.} Every halo has the pressure profile of
\citet{Arnaud2010}, which has the same shape as the finder's template. We also fix the hydrostatic mass bias
at $1-b_\mathrm{HSE}=0.8$, where $b_\mathrm{HSE}$ is the fraction by which the hydrostatic mass
underestimates the true mass. Because the filter is matched to one shape, we expect real clusters, whose
shapes differ from the template, to have lower and more scattered $q$. A catalogue model would have to
calibrate this effect, whereas a map-level simulator includes it once a range of profiles is painted. In
addition, a free mass bias is partly degenerate with $\asz$ in equation~(\ref{eq:y0}). Neither choice
changes the comparison between the forward models, but both would matter for real data.

\paragraph*{The 21-number summary.} We have kept the summary of \citet{Zubeldia2025}, four bins of $q$ by
five bins of redshift together with the mean lensing signal-to-noise, so that the forward models are
compared on the same data vector. On this summary, the map-level posterior is not more precise than the
catalogue-level ones (Fig.~\ref{fig:corner}), although the maps hold more information, in finer bins of
$q$, in the clusters below the threshold and in the tSZ field itself
\citep{Hill2014,Bolliet2018,EmbilVillagra2026}. Since the map-level simulator produces the maps, a richer
summary can be added without changing the forward model. We also expect a finer summary to make the
departures of catalogue models from the finder more visible, because the largest differences in the
counts are in the brightest bins (Fig.~\ref{fig:counts}).

\subsection{A possible offset in \texorpdfstring{$\omc$}{omega\_c}}\label{sec:discussion:omegac}

The map-level posterior is within 0.12 widths of the truth in every parameter at both survey sizes. Its most significant departure is in $\omc$, at $+0.10\pm0.07$ widths at 10 per cent of the sky and
$+0.10\pm0.04$ widths at 40.7 per cent (Table~\ref{tab:bias}). The second value differs from zero by 2.2
times its error, as computed from the unrounded values, but three points qualify this result. First, the two
measurements are not independent, because the mocks of the larger survey contain the skies of the smaller
one in their first 41 tiles. Second, all posteriors are applied to the same 100 mocks, so that a
fluctuation of the mock set would shift them together. The calibrated models also give $+0.09$ to
$+0.21$ widths in $\omc$. Third, at 10 per cent of the sky single networks range from $+0.02$ to $+0.16$
widths in $\omc$, and one network trained on 5000 simulations gives $-0.01$, comparable to this
seed-to-seed scatter (Appendix~\ref{app:robust}). At a tenth of a posterior width, the offset is small
compared with the slope bias of 0.5--0.7 widths left by the calibrated models, and none of our conclusions
depends on it. We have not determined whether it arises from the training or from the mocks, a question
that more mocks and more training simulations would settle.

\subsection{The road to real data}\label{sec:discussion:road}

In our simulations at goal depth, map-level SBI of the cluster counts is affordable at the scale of SO and
gives posteriors within 0.12 widths of the truth, with close to correct coverage. A catalogue model
calibrated on map-level mocks comes close at the true parameters, but it leaves a slope bias of 0.5--0.7 widths. Its coverage across the prior also depends on how the calibration follows the parameters
(Section~\ref{sec:discussion:calibrated}). The finder's response should therefore be forward-modelled in
an SO cluster-count analysis, at map level or through a calibration that follows the parameters and is
validated against map-level simulations. Before map-level SBI is applied to data, the simulator needs the
additions of Section~\ref{sec:discussion:limits}, namely clustered haloes, noise that follows the data,
foregrounds together with a spectrally constrained filter, and a model of the confirmation. It also needs
masks, a range of pressure profiles and a free mass bias. These steps each make a simulation more expensive, but the present cost (Table~\ref{tab:cost}) leaves room for them. The test used
here, in which catalogue-level and map-level posteriors are applied to the same map-level skies, can be
repeated at each step. This test measures directly whether a catalogue model is still good enough, and it
requires only map-level mocks and prior draws rather than a full training set.
